\renewcommand\runningtitle{Appendix A. Files on the Diskettes}
\chapter{Files on the Diskettes}

The two RetrOS/LT diskettes hold these files, in the order they are on
the diskettes. Install copies all of them to the RetrOS/LT directory
except those marked ``Not installed''; the two system files go to the
start of drive C, and \file{COMMAND.COM} to \file{C:\bs} as well. Each
diskette also has a hidden file, \file{RETROS.DSK}, by which Install knows
it. Every file is dated the day of the release.

{\small
\begin{longtable}{@{}P{1.05in}rP{0.85in}P{2.3in}@{}}
\toprule
\textbf{File} & \textbf{Bytes} & \textbf{Diskette} & \textbf{What it is}\\
\midrule
\endhead
\csname @@input\endcsname diskfiles.tex
\bottomrule
\end{longtable}}

\mwnewpart
\renewcommand\runningtitle{Appendix B. Technical Notes}
\chapter{Technical Notes}
\sectioncontents

\section{Version Numbers}\ixs{version, DOS}

\typed{VER} shows \shown{RetrOS/LT Version 1.0, DOS 2.10}. To a program
that asks through interrupt 21h, function 30h, RetrOS/LT reports DOS
version 2.10, as PC DOS 2.10 does. The commands supplied with RetrOS/LT
check for DOS 2.

\section{The System Files}\ixs{RETBIO.COM}\ixs{RETDOS.COM}

\file{RETBIO.COM} is RetrOS/LT's interface to the computer's BIOS: its
drivers for the keyboard and display, the printers, the serial port, the
clock and the disks. \file{RETDOS.COM} is the kernel, which gives programs
the DOS functions. Both are hidden. \file{RETBIO.COM} must be the first
entry in the root directory and be in one piece at the start of the data
area; \file{RETDOS.COM} must be the second entry. \typed{FORMAT /S} and
\typed{SYS} see to that.

The startup record of the disk loads \file{RETBIO.COM} at 0070:0000;
\file{RETBIO.COM} loads \file{RETDOS.COM}. On a fixed disk, the master
boot record that \typed{FDISK} writes first loads the startup record of
the active partition.

\section{Disks}\ixs{disks, formats}

\begin{tabularx}{\textwidth}{@{}P{1.0in}P{0.6in}P{0.6in}P{0.7in}P{0.6in}X@{}}
\toprule
\textbf{Format} & \textbf{Media byte} & \textbf{Sectors a cluster} & \textbf{Root entries} & \textbf{FAT sectors} & \textbf{Data starts at sector}\\
\midrule
360 KB & FD & 2 & 112 & 2 & 12\\
320 KB & FF & 2 & 112 & 1 & 10\\
180 KB & FC & 1 & 64 & 2 & 9\\
160 KB & FE & 1 & 64 & 1 & 7\\
Fixed disk & F8 & by size & 512 & by size & by size\\
\bottomrule
\end{tabularx}

A fixed disk's partition begins on head 1 of cylinder 0 and is described
by a 12-bit file allocation table. RetrOS/LT works out the layout of a
DOS partition from its size, as DOS 2 does: the smallest cluster that
keeps the number of clusters within 4,084, and as many table sectors as
that needs. Neither the PC nor the XT can tell when a diskette is changed,
so RetrOS/LT treats a diskette used in the last two seconds as unchanged,
and otherwise reads its file allocation table again when the disk's
buffers allow.

\section{Limits}\ixs{limits}

\begin{tabularx}{\textwidth}{@{}P{2.0in}X@{}}
\toprule
Fixed disks & The first two the BIOS reports, one DOS partition each.\\
Volume size & 65,535 sectors, about 32 MB.\\
File allocation tables & 12-bit.\\
Open files & 99, set by \typed{FILES=}; 20 for one program.\\
Path length & 63 characters from the root.\\
Command line & 127 characters.\\
Environment & 160 bytes to begin with.\\
Dates & 1980 to 2099.\\
\bottomrule
\end{tabularx}

\section{Building RetrOS/LT}

The kernel, \file{COMMAND.COM} and most commands are built from Microsoft's
MS-DOS 2 source, with a short series of changes applied in order, using
Microsoft's assembler and linker. \file{RETBIO.COM}, the startup records,
\typed{FORMAT}'s part for the IBM Personal Computer, the IBM commands,
\file{ANSI.SYS} and Install are new. The \textit{RetrOS/LT Guide to
Engineering} describes them, the changes and the build.

\mwnewpart
\renewcommand\runningtitle{Appendix C. EDLIN and DEBUG}
\chapter{EDLIN and DEBUG}
\sectioncontents

\section{EDLIN Commands}\ixs{EDLIN command}

\typed{EDLIN} works on numbered lines. In its commands a line may be given
as a number, as a period for the current line, as \typed{\#} for the line
after the last, or as \typed{+}\var{n} or \typed{-}\var{n} from the current
line. Type a line number alone to edit that line.

\begin{tabularx}{\textwidth}{@{}P{1.85in}X@{}}
\toprule
\textbf{Command} & \textbf{What it does}\\
\midrule
{}[\var{line}]\typed{I} & Inserts lines before \var{line}. End with \keys{Ctrl}{Break}.\\
{}[\var{line}][,\var{line}]\typed{L} & Lists lines.\\
{}[\var{line}][,\var{line}]\typed{P} & Lists lines a page at a time, and moves to the last.\\
{}[\var{line}][,\var{line}]\typed{D} & Deletes lines.\\
\var{line},\var{line},\var{line}\typed{M} & Moves the range to before the third line.\\
\var{line},\var{line},\var{line}[,\var{n}]\typed{C} & Copies the range, \var{n} times.\\
{}[\var{line}][,\var{line}][\typed{?}]\typed{S}\var{text} & Searches for \var{text}; \typed{?} asks at each.\\
{}[\var{line}][,\var{line}][\typed{?}]\typed{R}\var{old}\key{F6}\var{new} & Replaces \var{old} with \var{new}.\\
{}[\var{line}]\typed{T}\var{file} & Merges another file in before \var{line}.\\
{}[\var{n}]\typed{W}, [\var{n}]\typed{A} & Writes lines to disk, and appends more from it, for files too large for memory.\\
\typed{E} & Ends, saving the file; the old one is kept as \file{.BAK}.\\
\typed{Q} & Quits without saving, after asking.\\
\bottomrule
\end{tabularx}

\section{DEBUG Commands}\ixs{DEBUG command}

Numbers in \typed{DEBUG} are hexadecimal. An address is a segment and an
offset, \typed{1234:0100}, or an offset alone.

\begin{tabularx}{\textwidth}{@{}P{1.85in}X@{}}
\toprule
\textbf{Command} & \textbf{What it does}\\
\midrule
\typed{A} [\var{address}] & Assembles 8086 instructions as you type them.\\
\typed{C} \var{range} \var{address} & Compares two areas of memory.\\
\typed{D} [\var{range}] & Displays memory.\\
\typed{E} \var{address} [\var{list}] & Enters bytes into memory.\\
\typed{F} \var{range} \var{list} & Fills memory.\\
\typed{G} [=\var{address}] [\var{breakpoints}] & Runs the program.\\
\typed{H} \var{value} \var{value} & Shows the sum and difference.\\
\typed{I} \var{port} & Reads a port.\\
\typed{L} [\var{address} [\var{drive} \var{sector} \var{count}]] & Loads the named file, or disk sectors.\\
\typed{M} \var{range} \var{address} & Moves memory.\\
\typed{N} \var{file} & Names the file for \typed{L} and \typed{W}.\\
\typed{O} \var{port} \var{value} & Writes to a port.\\
\typed{Q} & Quits.\\
\typed{R} [\var{register}] & Shows or changes the registers.\\
\typed{S} \var{range} \var{list} & Searches memory.\\
\typed{T} [=\var{address}] [\var{count}] & Traces one instruction at a time.\\
\typed{U} [\var{range}] & Unassembles memory into instructions.\\
\typed{W} [\var{address} [\var{drive} \var{sector} \var{count}]] & Writes the file, or disk sectors.\\
\bottomrule
\end{tabularx}

\caution{\typed{W} with a drive and sector numbers writes straight to the
disk, past the file system. A wrong number destroys the disk's contents.}

\mwnewpart
\renewcommand\runningtitle{Appendix D. Credits and License}
\chapter{Credits and License}
\sectioncontents

\section{Where RetrOS/LT Comes From}\ixs{credits}\ixs{MS-DOS 2, derived from}

RetrOS/LT is a product of The Lyon Mountain Company, Republic Lyon
Mountain Corporation. It is derived from the source code of MS-DOS 2,
which was written by Microsoft Corporation and which Microsoft released to
the public in 2018 under the MIT License. The kernel, \file{COMMAND.COM},
\typed{SYSINIT}, and the commands \typed{CHKDSK}, \typed{DEBUG},
\typed{EDLIN}, \typed{EXE2BIN}, \typed{FC}, \typed{FIND}, \typed{MORE},
\typed{PRINT}, \typed{RECOVER}, \typed{SORT}, \typed{SYS} and most of
\typed{FORMAT} are that code; \typed{LINK} is Microsoft's program from the
same release.

RetrOS/LT is not MS-DOS and is not PC DOS, and neither Microsoft nor IBM
has any part in it. The Lyon Mountain Company added:

\begin{itemize}
\item \file{RETBIO.COM}, the startup records and \typed{FORMAT}'s part for
  the IBM Personal Computer, all new, and the kernel's console functions,
  which the released source lacks and which were reconstructed;
\item the commands written to behave as IBM's do: \typed{ASSIGN},
  \typed{BACKUP}, \typed{COMP}, \typed{DISKCOMP}, \typed{DISKCOPY},
  \typed{FDISK}, \typed{GRAPHICS}, \typed{MODE}, \typed{RESTORE},
  \typed{TREE} and \file{ANSI.SYS};
\item the Install program;
\item the RetrOS/LT name, banners and messages, the diskettes, and these
  guides.
\end{itemize}

IBM, IBM Personal Computer, PC DOS and Personal Computer XT are names of
International Business Machines Corporation; MS-DOS is a name of
Microsoft Corporation. They are used here only to say what RetrOS/LT runs
on and is compatible with.

\section{The Microsoft Notice}\ixs{MIT License}

The MIT License requires that this notice accompany the code. It is also
on the DOS diskette, as \file{NOTICE.TXT}.

\VerbatimInput[fontsize=\footnotesize,baselinestretch=1.0]{../../dist/NOTICE.TXT}

\section{The RetrOS/LT License}\ixs{license}\ixs{Single-Machine License}\ixs{Blanket License}

RetrOS/LT as The Lyon Mountain Company supplies it is licensed, not sold,
under the license that follows. It is on the DOS diskette as
\file{LICENSE.TXT}, and Install asks you to accept it.

\VerbatimInput[fontsize=\footnotesize,baselinestretch=1.0]{../../dist/LICENSE.TXT}
